% =====================================================================
% Kapitel 8 Fazit und Ausblick
% Studienarbeit: RAG vs. Fine-Tuning - Ein konzeptioneller Kriterienkatalog
% Autor: Tobias Weigold
%
% Zielumfang: ca. 1,5 Seiten
% Aufbau (wie in deinem Stub angelegt):
%   8.1 Zusammenfassung der Erkenntnisse
%   8.2 Beantwortung der Forschungsfragen
%   8.3 Ausblick
% =====================================================================

\section{Fazit und Ausblick}
\label{sec:fazit}

\subsection{Zusammenfassung der Erkenntnisse}

Die vorliegende Arbeit hat gezeigt, dass die Auswahl zwischen RAG, Fine-Tuning und hybriden Ansätzen für das betriebliche Wissensmanagement nicht auf eine rein technische Frage reduziert werden kann. In Kapitel~\ref{sec:kriterien} wurden fünfzehn Kriterien entlang der Dimensionen Wissen, Aufgabe, Ressourcen sowie Organisation und Regulatorik aus der Literatur abgeleitet und in Kapitel~\ref{sec:katalog} zu einem zweistufigen Entscheidungsmodell verdichtet. Regulatorische Anforderungen aus DSGVO und EU AI Act wirken dabei als K.-o.-Filter, während die verbleibenden dreizehn Kriterien in einer Ordinalmatrix bewertet werden. Multi-Hop-Anforderungen wurden als RAG-interne Verfeinerung modelliert, nicht als Auswahlkriterium zwischen den drei Optionen.

Inhaltlich hat sich RAG in dynamischen, nachweispflichtigen und regulierten Kontexten als strukturell vorteilhaft erwiesen. Fine-Tuning behält seine Berechtigung bei stilistisch und formatgetriebenen Aufgaben mit stabilem Wissensbestand, insbesondere bei Nutzung parameter-effizienter Verfahren wie LoRA \autocite[1]{Hu.2022}. Hybride Ansätze sind vor allem dort tragfähig, wo hohe Faktentreue und stark ausgeprägte Fachsprache zusammenfallen, sie erfordern aber eine empirische Vorprüfung an domänenspezifischen Referenzdaten \autocite[7]{Ovadia.2024}.

\subsection{Beantwortung der Forschungsfrage}

Die in Kapitel~\ref{sec:einleitung} formulierte Forschungsfrage nach den technischen, organisatorischen und juristischen Kriterien, anhand derer Unternehmen zwischen RAG und Fine-Tuning für ihr internes Wissensmanagement abwägen können, lässt sich damit dreiteilig beantworten.

Auf der \textit{technischen} Ebene sind Aktualisierungsfrequenz (\kref{1}), Wissensvolumen (\kref{2}), Strukturgrad (\kref{3}) und Faktentreue (\kref{4}) die zentralen Unterscheidungsmerkmale, wobei RAG bei hoher Dynamik, großem Volumen und starker Nachweispflicht strukturell überlegen ist. Für stilistisch und formatgetriebene Aufgaben (\kref{5}) sowie latenzkritische Einsätze (\kref{11}) bleibt Fine-Tuning die belastbare Option. Auf der \textit{organisatorisch-wirtschaftlichen} Ebene sind Rechenressourcen (\kref{8}), Datenaufwand (\kref{9}), Betriebskosten (\kref{10}), Verfügbarkeit von Fachexpertise (\kref{14}) und Vendor Lock-in (\kref{15}) die maßgeblichen Kriterien. Sie sprechen mehrheitlich für RAG, sofern nicht ein spezifischer Aufgabenzuschnitt Fine-Tuning nahelegt. Auf der \textit{juristischen} Ebene führen die Anforderungen des EU AI Act an Protokollierung und Transparenz \autocite[Art.~12, Art.~13]{EuropaischesParlamentundRat.2024} sowie das Recht auf Löschung nach Artikel~17 DSGVO \autocite[Art.~17]{EuropaischesParlamentundRat.2016} in regulierten Anwendungsfällen zu einem strukturellen Vorteil von RAG. Fine-Tuning eines GPAI-Modells löst zudem eine potenzielle Anbieterrolle nach Artikel~25 Absatz~1 lit.~b aus, sobald eine substantielle Modifikation vorliegt \autocite[Rn.~61]{EuropeanCommission.2025}. Damit erweitert die Arbeit die bisher primär technisch geführte Diskussion um eine explizit regulatorische Dimension.

\subsection{Ausblick}

Aus den Ergebnissen ergeben sich drei zusammenhängende Forschungsperspektiven. Erstens ist die entwickelte Matrix empirisch zu validieren. Sinnvoll erscheint eine retrospektive Anwendung auf abgeschlossene Enterprise-Projekte in Kombination mit strukturierten Experteninterviews, um die Ordinalbewertungen und die Gewichtung der Kriterien an konkretem Erfahrungswissen zu spiegeln.

Zweitens verändern sich die technischen Grundlagen des Vergleichs weiter. Agentische RAG-Architekturen, die Retrieval, Planung und Werkzeugnutzung in mehrstufigen Zyklen verschränken, sowie kontinuierliches Lernen auf Modellseite verschieben die Grenzen zwischen den in dieser Arbeit unterschiedenen Optionen. Auch die stärkere Integration von Wissensgraphen in den Retrieval-Prozess \autocite[17]{Karakurt.2026} eröffnet neue hybride Konfigurationen, die im vorliegenden Kriterienkatalog nur als Verfeinerung von K3 und K6 angelegt sind und in einer Folgearbeit als eigenständige Option modelliert werden könnten.

Drittens ist die regulatorische Konkretisierung des EU AI Act durch Leitlinien und Durchführungsakte fortlaufend zu beobachten. Insbesondere die praktische Ausgestaltung der Schwellenwerte für substantielle Modifikationen \autocite[Rn.~61]{EuropeanCommission.2025} und die branchenspezifischen Umsetzungsanforderungen für Hochrisiko-Anwendungen werden die relative Attraktivität von Fine-Tuning gegenüber RAG in den kommenden Jahren spürbar prägen. Die vorgeschlagene Matrix ist so konstruiert, dass sich derartige Änderungen ohne strukturellen Umbau in aktualisierten Bewertungen abbilden lassen. Damit bleibt der Kriterienkatalog auch unter fortschreitender technischer und regulatorischer Entwicklung ein tragfähiger Bezugsrahmen für die architektonische Auswahlentscheidung im Enterprise-Wissensmanagement.
